% SGA 3, 第 VI_B 篇, 第 1 轮英文重构.
% 范围: 引理 5.6.1 全文, 含底本编者注 41.
% 底本: Polo--Gille Exp6B-13oct24.pdf, 分篇第 29--31 页
% (合订本第 370--372 页; 原版页边印刷页码 353--354, 重版本页眉 359--361).
% 引理 5.6.2.0 始于分篇第 31 页 / 合订本第 372 页.
% OCR 与 Reinhold 英文 Markdown 仅为比对材料.
% 英文底本旧注: 四图均为第 1 轮的 180 dpi 底本精确裁切; 下方现为原生 TikZ 图码.

\medskip
\phantomsection\label{lem:sga3-VIB-5.6.1}%
\noindent\textbf{引理 5.6.1.---}
\emph{设 \(S\) 是概形, \(f:X\to S\) 是局部有限呈示的 \(S\)-概形,
\(s\) 是 \(S\) 的点, \(x\) 是 \(X_s\) 的闭点.
假设 \(f\) 在 \(X_s\) 的某个包含 \(x\) 的不可约分支
\(Z\) 的泛点 \(z\) 处普遍开, 且}
\[
  \dim_x(X_s)=\dim Z.
\]
\emph{则存在交换图}
\begin{center}
  \phantomsection\label{diag:sga3-VIB-5.6.1-statement}%
  \begin{tikzcd}[row sep=large,column sep=large]
    X \arrow[r,"f"] &
    S \\
    S''
      \arrow[u,"h"]
      \arrow[r,"\pi"]
      \arrow[ur,"\phi"] &
    S' \arrow[u,"w"']
  \end{tikzcd}\enspace,
\end{center}
\emph{其中 \(S'\) 仿射, \(w\) 平展,
\(\pi\) 有限、满射且有限呈示, 而 \(\phi^{-1}(s)\)
仅由一点 \(s''\) 组成, 使得 \(h(s'')=x\) 且 \(\kappa(s'')=\kappa(x)\).}%
\footnote{\label{note:sga3-VIB-41}编者注: 因此,
\(h\) 诱导 \(X\times_S S''\to S''\) 的一个截面 \(\sigma\),
使 \(\sigma(s'')\) 位于 \(x\) 上; 比较 EGA~IV\(_3\), 14.5.10.}

\emph{证明.}
首先可假设 \(S=\operatorname{Spec}(A)\) 及 \(X=\operatorname{Spec}(B)\),
其中 \(B\) 是有限呈示 \(A\)-代数.
设 \(\mathfrak p\) 和 \(\mathfrak q\) 分别是 \(A\) 和 \(B\)
中对应于 \(s\) 和 \(x\) 的素理想, 从而
\[
  \mathfrak p=A\cap\mathfrak q.
\]
置 \(n=\dim_x(X_s)\), 并选取 \(t_1,\ldots,t_n\in B\),
使其在
\[
  \mathcal O_{X_s,x}=B_{\mathfrak q}/\mathfrak p B_{\mathfrak q}
\]
中的像构成一组参数系. 则
\[
  \mathcal O_{X_s,x}/(t_1,\ldots,t_n)
\]
在 \(\kappa(x)\) 上有限维,
因而在 \(\kappa(s)\) 上也有限维,
因为 \(x\) 是代数 \(\kappa(s)\)-概形 \(X_s\) 的闭点.
所以 \(S\)-态射
\[
  u:X\longrightarrow\mathbb A^n_S
    =\operatorname{Spec}\bigl(A[T_1,\ldots,T_n]\bigr),
  \qquad T_i\longmapsto t_i,
\]
有限呈示, \(x\) 在纤维 \(u^{-1}(u(x))\) 中孤立, 且
\[
  u(x)=\tau_0(s),
\]
其中 \(\tau_0:S\to\mathbb A^n_S\) 表示 \(\mathbb A^n_S\to S\)
的零截面, 对应于 \(A\)-代数同态
\[
  A[T_1,\ldots,T_n]\longrightarrow A,\qquad T_i\longmapsto0.
\]
\(X\) 中在其位于 \(\mathbb A^n_S\) 上的纤维中孤立的点所成的集合开
(EGA~IV\(_3\), 13.1.4). 因此缩小 \(X\) 后,
可假设 \(u\) 拟有限, 且 \(u^{-1}(u(x))=\{x\}\).
由\hyperref[lem:sga3-VIB-5.6.1.0]{引理 5.6.1.0}, \(u\) 在 \(x\) 处普遍开.

设
\[
  B_0=B/(t_1,\ldots,t_n),\qquad
  X_0=X\times_{\mathbb A^n_S}\tau_0(S)
      =\operatorname{Spec}(B_0),
\]
并设 \(u_0:X_0\to S\) 是 \(u\) 作基变换
\(\tau_0:S\hookrightarrow\mathbb A^n_S\) 所得的态射.
则 \(u_0\) 拟有限且有限呈示, 在 \(x\) 处普遍开,
且 \(x\) 是 \(X_0\) 中位于 \(s\) 上的唯一点.

设 \(A'\) 是局部环
\[
  A_{\mathfrak p}=\mathcal O_{S,s},
\]
的 Hensel 化, 并置
\[
  S'=\operatorname{Spec}(A'),\qquad
  B'_0=B_0\otimes_A A',\qquad
  X'_0=X_0\times_S S'=\operatorname{Spec}(B'_0).
\]
\(S'\) 的闭点 \(s'\) 是 \(S'\) 中位于 \(s\) 上的唯一点,
且 \(\kappa(s')=\kappa(s)\).
概形 \(X'_0\) 在 \(s'\) 上有唯一点 \(x'\);
此外 \(\kappa(x')=\kappa(x)\), 且 \(x'\) 也是 \(X'_0\)
中同时位于 \(s\) 和 \(x\) 上的唯一点.
由于 \(A'\) 是 Hensel 环, EGA~IV\(_4\), 18.5.11 给出既开又闭部分的不交分解
\[
  \tag{*}\label{sga3:vi:lemma:5-6-1:equation:star}
  X'_0=V\amalg W,
  \qquad
  V=\operatorname{Spec}(\mathcal O_{X'_0,x'}).
\]
局部环 \(\mathcal O_{X'_0,x'}\) 在 \(A'\) 上有限且有限呈示.
因此 \(u'_0\) 在 \(V\) 上的限制 \(\pi\) 有限且有限呈示.
此外, 由于 \(u'_0:X'_0\to S'\) 在 \(x'\) 处开,
\[
  u'_0(V)=\operatorname{Spec}(\mathcal O_{S',s'})=S',
\]
所以 \(\pi\) 满射.

这证明了 \(S=S'\) 时的断言. 一般地,
\(A'\) 是各个 \(A\)-子代数 \(A_i\) 的滤过正向极限,
这些子代数在 \(A\) 上平展, 且
\[
  S_i=\operatorname{Spec}(A_i)
\]
在 \(s\) 上有唯一点 \(s_i\), 满足 \(\kappa(s_i)=\kappa(s)\). 有
\[
  B'_0=\varinjlim_i B_i,
  \qquad
  B_i=B_0\otimes_A A_i.
\]
置
\[
  X_i=\operatorname{Spec}(B_i)=X_0\times_S S_i,
  \qquad
  C=\mathcal O_{X'_0,x'}.
\]
由 \((*)\), 对某个幂等元 \(f\in B'_0\), 有
\[
  C\simeq B'_0/fB'_0
\]
存在一个指标 \(i\) 和一个在 \(B'_0\) 中的像为 \(f\) 的元素 \(f_i\in B_i\).
置
\[
  C_i=B_i/f_iB_i,\qquad V_i=\operatorname{Spec}(C_i).
\]
则
\[
  C=C_i\otimes_{A_i}A',
  \qquad
  V=V_i\times_{S_i}S',
\]
且 \(C_i\) 是有限呈示 \(A_i\)-代数, 因为 \(B_i\) 也是如此.
所以态射
\begin{center}
  \phantomsection\label{diag:sga3-VIB-5.6.1-limit-source}%
  \begin{tikzcd}[row sep=large,column sep=large]
    X'_0 \arrow[dr,"u'_0"] &
    {} \\
    V
      \arrow[u,"h'"]
      \arrow[r,"\pi"] &
    S'
  \end{tikzcd}
\end{center}
由有限呈示态射
\begin{center}
  \phantomsection\label{diag:sga3-VIB-5.6.1-limit-descent}%
  \begin{tikzcd}[row sep=large,column sep=large]
    X_i \arrow[dr,"u_i"] &
    {} \\
    V_i
      \arrow[u,"h_i"]
      \arrow[r,"\pi_i"] &
    S_i
  \end{tikzcd}
\end{center}
作基变换 \(S'\to S_i\) 得到.

对每个 \(j\geq i\), 令
\[
  V_j=V_i\times_{S_i}S_j
\]
并设 \(\pi_j:V_j\to S_j\) 是 \(\pi_i\) 作基变换所得的有限呈示态射.
由于 \(\pi:V\to S'\) 有限且满射,
EGA~IV\(_2\), 8.10.5 表明, 对某个 \(j\),
态射 \(\pi_j:V_j\to S_j\) 有限且满射.
则 \(w:S_j\to S\) 仿射且平展,
\(S_j\) 在 \(s\) 上有唯一点 \(s_j\),
而 \(V_j\) 在 \(s_j\) 上有唯一点 \(x_j\),
因为 \(x'\) 是
\[
  V=V_j\times_{S_j}S'
\]
中位于 \(s'\) 上的唯一点:
\begin{center}
  \phantomsection\label{diag:sga3-VIB-5.6.1-final-composite}%
  {\small
  \begin{tikzcd}[row sep=large,column sep=large]
    X_j
      \arrow[r]
      \arrow[dr,"u_j"] &
    X_0
      \arrow[r]
      \arrow[dr,"u_0"] &
    X
      \arrow[d,"f"']
      \arrow[dr,"u"] &
    {} \\
    V_j
      \arrow[u,"h_j"]
      \arrow[r,"\pi_j"] &
    S_j
      \arrow[r,"w"] &
    S
      \arrow[r,"\tau_0"] &
    \mathbb A^n_S
  \end{tikzcd}}
\end{center}
所以 \(x_j\) 是 \(V_j\) 中位于 \(s\) 上的唯一点,
它在 \(V_j\to X_j\to X\) 下的像为 \(x\), 且
\[
  \kappa(x_j)=\kappa(x).
\]

% 翻译覆盖：本组件全文；源文件 43_expose_VIB_lemma561_loop2_r1_en.tex。
% 自检：行内数学 127，陈列数学 21，原生 tikzcd 4；全部数学及编号保留，公式内自然语言已译。
% 引用键核对：label 7，hyperref 1，SGARefLink 0，行内目标 0；脚注 1，星号脚注正文 0，均与源文件一致。
% 编译及渲染审查由主线程整合完成；此处未宣称编译或独立校对完成。
